\documentclass[12pt, letterpaper]{article}
\usepackage[letterpaper,portrait,top=0.4in, bottom=0.4in, left=0.6in, right=0.6in]{geometry}

\usepackage{amsmath, amsfonts, amsthm, amssymb}
\usepackage{graphicx, float}
\usepackage{mathtools}
\usepackage{titlesec}
\usepackage{interval}
\usepackage{hyperref}
\usepackage{siunitx}
\usepackage{titling}
\usepackage{vwcol}
\usepackage{setspace}
\usepackage{empheq}
\usepackage{cancel}
\usepackage{esdiff}
\usepackage{multicol}
\usepackage{graphicx}
\usepackage{mdframed}
\usepackage{esdiff}
\usepackage{tikzsymbols}
\usepackage{multicol}
\usepackage{tikz}
\usepackage{varwidth}
\usepackage{pgfplots}
\pagestyle{empty}
\pgfplotsset{compat=1.18}
\intervalconfig {
	soft open fences
}
\newcommand\HUGE{\@setfontsize\Huge{50}{50}}
\titleformat{\section}
{\normalfont\large\bfseries}{Article~\thesection.~}{0em}{}

\begin{document}
\setstretch{1.5}
\fontsize{12pt}{12pt}\selectfont
\setlength{\abovedisplayskip}{0pt}

\begin{center}
	\vspace*{\fill}

	\vspace{-80pt}
	\huge
	\textsc{Constitution of the}

	\vspace{2pt}
	\huge
	\scalebox{1.7}{\fontsize{24pt}{0pt}\selectfont \textsc{Peddie Math Club}}

	\vspace{40pt}
	\includegraphics[scale=3]{logo.png}

	\vspace*{\fill}
\end{center}
\newpage

\section*{Preamble}
Taking as point of departure the historic fact that the mathematicians, students and scientists of all the grades and ages of The Peddie School, united in the Club of Mathematics headed by our honorable President Jiayu “Jerry” Liu, by their struggle in the formation of the Club and Mathematical Revelation overthrew the old class order based on ageism, sleep deprivation and the fear of a lower GPA and began to build a society in which human intellect and man would be emancipated from exploitation and arbitrariness, a society in which each student and all of them together would find conditions for free and all-round development; bearing in mind that the development of the person’s intellectual foundations and of mathematical-social relationships, and the further development of relations based on self-studying and mathematical equality, have brought about substantial changes in scholarly and social relations which call for corresponding changes in the Constitution of the Math Club of 2022, but that in which the Constitution in individual spheres has not yet been changed by constitutional amendments; wishing to consolidate and further develop the revolutionary achievements already accomplished, to strengthen the right and responsibility of the Members of the Club for their own development and for the development of the Peddie community as a mathematical whole, and to ensure the further development of mathematical-democratic relations of the inherent interest in mathematics on the road to the emancipation of thought and the creation of a mathematical society; proceeding also from the need to strengthen the constitutional system on uniform foundations of mathematical acquisition of knowledge by the new Constitution of the Club of Mathematics, which, in addition to the changes it introduces, incorporates and adjusts to it the provisions of the 2022 Constitution of the Club of Mathematics. The Officers, in agreement with the Students and Faculty of the Department of Mathematics, have adopted \textsc{\large The Constitution of the Math Club}.

\section{Mission}
Our core mission is to propagate the passion for mathematics to everyone at Peddie. Regardless of their mathematical prowess, everyone can participate and engage in establishing an inclusive community of thinkers. We will create a comfortable space for avid mathematicians to express their passion that will emanate to other students at Peddie. We will offer opportunities to Peddie Students in a broader scope not restricted to the classroom and the traditional math curriculum, paving the way to develop intellectual talents capable of altering the globe. 

\section{Officers}
There shall be three standing club officers: the President, the Vice President and the Secretary. Additional officers may be appointed as outlined under Article 6.

\section{Appointment of the President}
At the final meeting of the Club in the current school year, an Electoral College shall be convened with the President, the Vice-President and the Secretary as its \textit{ex officio} members. Each member of the College shall nominate one candidate for the office of President. If one person is nominated by two or more members, they shall be appointed President of the Math Club for the following school year. Should no candidate receive a nomination from two or more members, further discussion and nominations shall occur until a new President is agreed upon and appointed.

\section{Appointment of the Vice-President and Secretary}
On the appointment of the President for the following year, the newly appointed President shall appoint the following year’s Vice-President and Secretary.

\section{Responsibilities of the President and Vice-President}
The President and Vice-President shall organize and prepare future activities of the Club and are responsible for delivering said objectives on a timely basis with respect to the appointed date and time decided by the members of the Club. The Vice President shall ensure that the Secretary does not engage in activities deemed to be revolutionary in nature. The President holds to power to reconsider any major decisions made by the Vice-President and the Secretary pursuant to the following procedure:
\begin{enumerate}
	\item Hold a meeting consisting of the three standing officers of the Club to discuss whether the decision is unconstitutional or unethical.
	\item If the disagreement persists, the Chair of the Department of Mathematics shall settle the dispute.
\end{enumerate}

\section{Responsibilities of the Secretary}
The Secretary shall be responsible for maintaining the veracity and timeliness of the contents of the Google Document entitled ``TODO'' with regard to the content discussed at every meeting of the Club. The Secretary shall also appoint, with the advice and consent of one or both of the President and Vice-President, one or more individuals to discharge the duties of communication with regard to the following people and organizations, as well as the function and purpose for the communication with the person in question:
\begin{enumerate}
	\item The faculty with the responsibilities that is held by Mr. Gregory Wriede in the ninth month of the two-thousandth and twenty-fourth year of the Common Era, who is responsible for the distribution of club funding.
	\item Those who have graduated from The Peddie School in periods of time, at the time of communication, which are regarded as having occurred in the past, and who may be connected to this Club.
	\item The Chair of the Department of Mathematics at The Peddie School, who may also be an advisor.
	\item The faculty with the responsibility for the coordination and effective organization of mathematical contests and competitions, a responsibility that is held, at the time of writing, by Mrs. Victoria Montgomery.
\end{enumerate}
The Secretary shall also appoint an individual to maintain the Club website. In times when a burden arises with no obvious individual from the above list to address, the Secretary may temporarily appoint a person responsible for said unhandled business.

\section{Dismissal of Officers}
Each Officer of the Club serves at the pleasure of the Chair of the Department of Mathematics, and may be dismissed from office when they order their removal. Each Officer may also be dismissed following a vote by the other two Officers of the Club. The dismissal of the President must be approved by the Chair of the Department of Mathematics. 

\section{Resignation}
Each Officer of the Club may resign from their position for any reason at any time.

\section{Appointments Following the Removal or Resignation of Officers}
An Officer is removed from their position in case of a Dismissal pursuant to Article 7, their expulsion from The Peddie School, or are otherwise unable to discharge the duties of their office. In the event of the removal of the Vice-President or Secretary, the President shall appoint their replacement. In the event of the removal of the President, the Chair of the Department of Mathematics shall appoint a replacement from a list of two candidates proposed by the Vice-President and Secretary.

\section{Amendments to the Constitution}
Amendments may be proposed by any member of the club and adopted if  two-thirds or more members cast a vote in favor.

\section{Licensing of this Constitution}
This work is licensed under the Creative Commons Attribution-ShareAlike 4.0 International License. To view a copy of this license, visit \underline{\href{https://creativecommons.org/licenses/by-sa/4.0/}{\color{blue}{https://creativecommons.org/licenses/by-sa/4.0/}}}.

\end{document}
